\subsection{Program Loading and Execution Mechanics}

\subsubsection{The Lifecycle of Transformation: From Storage Binary Block to Active Virtual Memory Process}

A compiled program on disk is a formatted artifact, not a running computation. Execution begins when the operating system accepts an executable, constructs a process container, maps loadable regions into a virtual address space, prepares stack and startup metadata, and transfers control to an entry point. The chain is: \textbf{stored binary} $\rightarrow$ \textbf{validated image description} $\rightarrow$ \textbf{mapped virtual memory} $\rightarrow$ \textbf{initial execution context} $\rightarrow$ \textbf{running process}.

In the bare-metal C world, this is like surveying a cold city grid: the loader places fixed districts (text, data, BSS) at predetermined addresses, opens a construction lane for the stack, and hands the CPU a starting street coordinate. Nothing is ``Python-aware'' yet; the process image is native machine code plus static data regions.

\subsubsection{The OS Loader Subsystem: VMA Memory Mapping, Page Table Initialization, and Stack Argument Injection}

The loader's job is to translate on-disk segment descriptions into live mappings. Code pages are executable and read-only; data pages are writable; constants may be read-only. The stack and heap begin as anonymous regions grown on demand. Command-line arguments, environment blocks, and dynamic-linker metadata are injected so startup code can find \texttt{argc}/\texttt{argv} and dependent shared libraries.

Conceptually, think of the OS as assigning apartment numbers in a building the process believes it owns alone. Virtual addresses are the door numbers; page tables are the building directory maintained by the kernel. A C program's pointers are interpreted inside that private numbering scheme.

\subsubsection{Relocation Invariants and Control Handover to the Hardware Thread Execution Context}

Linked executables still carry relocation fixups: symbol addresses that must be patched to match where segments actually mapped. Once relocations are applied and dependencies resolved, the kernel sets the instruction pointer and stack pointer according to the platform ABI and begins native execution. From this moment, user instructions run directly on the CPU until a syscall, fault, interrupt, or trap returns control to the kernel.

\subsubsection{Native Entry Points: Loader Handover, C Runtime Startup, and the Call to \texttt{main()}}

The visible \texttt{main()} is not the first instruction executed. Startup code initializes runtime services (TLS, constructors, stdio hooks), then calls \texttt{main()}. When \texttt{main()} returns, teardown paths run and the process exits. This entire story is the \textbf{Grid}: rigid districts, hardware stack discipline, and a single native call chain anchored in machine frames.

Python's launch differs only at the front door: the OS loads \texttt{python}, not the \texttt{.py} file as a native executable. The user's script is input to an already-hosted runtime that will create its own internal ``hotel'' inside the process heap, described in the next sections.
